% =========================================================
%  frontmatter/acknowledgements.tex
%  Ucapan terima kasih
% =========================================================

\chapter*{Ucapan Terima Kasih}
\addcontentsline{toc}{chapter}{Ucapan Terima Kasih}
\thispagestyle{plain}

Buku ajar ini tidak akan selesai tanpa dukungan berbagai pihak. Penulis
menyampaikan terima kasih kepada pimpinan \NamaPT, Dekan \NamaFakultas,
dan Ketua Program Studi \NamaProdi\ atas kesempatan dan dukungan yang
diberikan dalam penyusunan bahan ajar ini.

Terima kasih juga disampaikan kepada rekan dosen di lingkungan Program
Studi \NamaProdi\ yang telah memberikan masukan atas struktur materi,
khususnya pada bagian asesmen berbasis capaian pembelajaran. Masukan
tersebut sangat membantu memastikan bahwa setiap bab benar-benar terhubung
dengan Sub-CPMK, bukan hanya memuat materi yang menarik.

Penulis juga menghargai para penulis dan pengembang yang karyanya menjadi
sumber utama buku ini, yaitu dokumentasi resmi Python, Pygame, dan Unity,
buku \istilah{Making Games with Python \& Pygame} serta \istilah{Invent Your
Own Computer Games with Python} karya Al Sweigart, tulisan Robert Nystrom
tentang \istilah{Game Loop}, penjelasan Glenn Fiedler mengenai
\istilah{timestep}, bahan kuliah Edirlei dos Santos dan Rinaldi Munir
mengenai kecerdasan buatan permainan serta algoritma A*. Karya-karya itu
membuat penulisan buku ajar ini menjadi mungkin.

Akhirnya, terima kasih kepada mahasiswa yang telah dan akan mengikuti mata
kuliah ini. Pertanyaan, kesalahan yang ditemukan, dan kesulitan yang
dihadapi ketika mengerjakan praktikum adalah bahan perbaikan yang paling
jujur bagi buku ini.

\vspace{0.8cm}

\begin{flushright}
\BookAuthor
\end{flushright}

\clearpage
